\subsection{What the Host estimates leave open}
\label{sec:host-uniformity-gate}

\paragraph{K17-HOST-GREENBERG-GATE.}
The established Step~17, Host and window estimates do not by themselves
remove the hypothesis that the one-variable plus Iwasawa module is
finite.  The following argument checks the full quantifiers and gives
explicit \emph{compatibility countermodels}: abstract modules and
artificial trace polynomials satisfying the numerical properties under
discussion.  They are not arithmetic counterexamples, actual
$K$-groups, or actual automorphic trace systems.  The conclusion
concerns the proposed inference from these estimates; it does not
exclude a different use of the stipulated arithmetic source contract.

\paragraph{The constants and the tower parameter.}
The source's \texttt{CountI} first fixes $\chi$, then chooses
$C,r_1,r_2$, and bounds the indicated sum of normalized lengths by
$(Nt)^C$ over its admissible hosts and Case~I primes.
The full \texttt{PointHV} has the stronger but still ordered form
\[
 \begin{gathered}
 \forall\chi\ \forall j:\mathbb N\longrightarrow\mathbb N\
 \exists C,r_1,r_2\ \forall h,t,B:\\
 \bigwedge_{\ell\in\mathcal B_\chi}
     \operatorname{LayerLe}(\chi,\ell,j(\ell),t,B)\\
 \Longrightarrow\quad
 \sum_{\ell}\log\ell\sum_{1\le e<t}
       \operatorname{length}\mathcal S_\ell(E_h,e)^{+_e}
       \le (Nt)^C+CB.
 \end{gathered}
 \tag{HU.1}
\]
All sums over primes in (HU.1) mean the source's arbitrary finite sets
of odd primes.  The admissibility conditions and $B\ge0$ are as in the preceding
Host theorem.  For each fixed $\ell$, $j(\ell)$ is one fixed
cyclotomic layer depth.  Replacing it by larger depths permits the
constants to change.  Moreover $B$ is an input estimate for the
finite-layer groups, and the output retains it; it is not an
independently obtained uniform bound on the tower.

More explicitly, \texttt{LayerLe} ranges over presentations of
$\mathbb Q(\mu_{\ell m})$ and cyclotomic $\ell$-power extensions of
degree $\ell^{j(\ell)}$.  The definition \texttt{IsLayer} requires
both that cyclotomic condition and that degree.  Its conclusion
\texttt{VVat} bounds the sum of strict Selmer lengths for $1\le e<t$.
The exact source locators are
\path{HeightAndVanishing/Counts.lean:CountI},
\path{HeightAndVanishing/LayerLe.lean:LayerLe,PointHV},
\texttt{StmtFields.lean:IsLayer}, and
\texttt{LengthEqAddEigen.lean:VVat}, relative to
\path{src/Development/AutomorphicArithmetic/AutomorphicVarieties/}
in the pinned source \cite{K17PinnedLeanMirror}.
The final combination is
\path{HeightAndVanishing/PointArith.lean:point}.

There is also a field-theoretic restriction on this uniformity.
For a fixed character of modulus $m$ and order $o$, the host is
\[
 E_h=F_\chi\,\mathbb Q(\mu_{2^{a+2}})^+
           \mathbb Q(\sqrt{-q})\,\mathbb Q(\sqrt{-d}),
 \qquad N=\operatorname{lcm}(2^{a+2},m)qd.
 \tag{HU.2}
\]
Every field adjoined to $F_\chi$ in (HU.2) has degree a power of
two.  Hence $[E_h:F_\chi]$ is a power of two.  Since the fields
are abelian over $\mathbb Q$, for each odd prime $p$ the restriction
map identifies the $p$-primary parts of
$\operatorname{Gal}(E_h/\mathbb Q)$ and
$\operatorname{Gal}(F_\chi/\mathbb Q)$.  Their order is the fixed
number $p^{v_p(o)}$, independent of $a,q,d$.
Consequently this host family does not contain an unbounded
$p$-cyclotomic tower.  Such a specialization changes $\chi$ and
therefore permits the constants to change.  When $p$ divides the
cyclotomic modulus, it also lies in $\mathcal B_\chi$ and its layer
estimate occurs explicitly in (HU.1).

Restriction and norm through the extension in (HU.2) multiply a
$p$-primary group by a $p$-adic unit.  Whenever the local conditions
are preserved, this transports an upper bound to a direct summand;
it does not make the bound decrease.  Indeed the host condition
$q,d>2t+1$ forces $N>4t^2$.
Even a bound polynomial in conductor with universal exponents
would allow positive Iwasawa $\lambda$: at conductor $p^{n+1}$
it is exponential in $n$, whereas the model below grows only
linearly in $n$.

\paragraph{The downstream window theorem.}
The checked endpoint has the definition
\[
 \texttt{VWin}:\quad
       \forall\chi,\ \texttt{VWinAt}(\chi,401).
\]
Its source is \path{CharacterBoxes/ExistsPageN.lean}.
Here the directory prefix is
\path{src/Development/AutomorphicArithmetic/}.
For each fixed $\chi$, each $0<\varepsilon<1/52$, and each function
$a(x)$ with
\[
 x^{2+8\varepsilon}\le 2^{a(x)}
      \le 2x^{2+8\varepsilon},
\]
it chooses $J\ge1$ and $X_0$.  For every $X\ge X_0$ and every
pair $(q,d)$ satisfying \texttt{WinPair} at
$(\varepsilon,X^{1/J})$, it then chooses $Y$ with
$X^{1/J}\le Y$ and $Y^2\le X$.  For every $Y\le x\le Y^2$
and every field satisfying \texttt{IsChiField}, its conclusion is
\[
 \sum_{\ell}\log\ell
     \sum_{1\le e<\lfloor x^{401}\rfloor}
           \operatorname{length}\mathcal S_\ell(E,e)
      \le \sqrt{\exp(x^\varepsilon)}
      =\exp(x^\varepsilon/2).
 \tag{HU.3}
\]
This includes both eigenspaces and the prime two.  The proof is
\path{CMFieldExponentialBounds/AdmissibleOfPageN.lean},
declarations \texttt{vwinAt\_of\_joint} and
\texttt{vwin\_of\_joint}; its complete argument and the
complete endpoint definitions were read and their retrieved source
blobs verified at the pinned commit \cite{K17PinnedLeanMirror}.
The intermediate definitions \texttt{ThmBab,ThmHV,JointWin} are in
\texttt{PrimeSetBoundDataOfPrimitiveChar.lean}.
The full proof \path{HeightAndVanishing/ThmHV.lean:thmHV}
chooses the fixed depth
$j(\ell)=\lfloor\log_\ell R_\ell\rfloor+3$ before applying
\texttt{PointHV}.

\begin{sloppypar}
The endpoint has a closed source composition under the same stipulated
contract as the preceding Host result.  Apply
\path{UpperTheory.V.v_window_of_literal_automorphy_n}
to \path{UpperTheory.V.SmallClassNumbers.litN}.
The latter is in \path{SmallClassNumbers/CorGKOf.lean}; the former
is the final theorem of \path{WindowBoundOfLiteratureLevel.lean}.
Its directory is
\path{ModularityArgument/ThresholdArguments/}\allowbreak%
\path{DenseThresholdCases/}.
For the two premises in \texttt{vwin\_of\_joint}, the intervening
\texttt{HVCore.lean:of\_four\_alt}, in that same directory,
uses two explicit providers.

First, \texttt{Win.thmBab\_of\_lemmaK} gives \texttt{ThmBab}
using Poitou--Tate, the local inputs, Ferrero--Washington and
the established \texttt{LemmaK} witness.
Second, \texttt{Odd.thmBcFree\_of} gives \texttt{ThmBcFree}
when applied to \texttt{ML.iv} with the Poitou--Tate and unit inputs.
The complete body of \texttt{of\_four\_alt} was read.
Then \path{HeightAndVanishing/JointWinOfCounts.lean:vwin_of_counts}
uses \texttt{jointWin\_of\_counts} to produce
$\texttt{JointWin}(401)$ and \texttt{thmBc\_iff} to convert
\texttt{ThmBcFree} into $\texttt{ThmBc}(401)$.
Thus the conditional premises of \texttt{vwin\_of\_joint} are
supplied by explicit source declarations; they are not omitted
assumptions.  The compatibility argument below would also apply
if (HU.3) were merely imposed as a stronger numerical test.
\end{sloppypar}

For every fixed $\varepsilon>0$, every fixed polynomial in $x$
is eventually smaller than $\exp(x^\varepsilon/2)$: taking
logarithms reduces this to $\log x=o(x^\varepsilon)$.
Thus (HU.3) permits such polynomial contributions even without
selecting favorable windows.  Letting $\varepsilon$ tend to zero
at a fixed $x$ would reverse the quantifiers: both $J$ and $X_0$,
as well as the prescribed hosts, may change.

\paragraph{An explicit nonfinite module.}
Let $p$ be odd, let $\Gamma\cong\mathbb Z_p$ have generator $\gamma$,
put $u=1+p$, and set
\[
 \Lambda=\mathbb Z_p[[\Gamma]],\qquad
 M=\Lambda/(\gamma-u).
 \tag{HU.4}
\]
With $\gamma=1+T$, this is
$\mathbb Z_p[[T]]/(T-p)\cong\mathbb Z_p$.
It is infinite, $\Lambda$-torsion and $\mathbb Z_p$-torsion free,
with $\mu=0$ and $\lambda=1$.
Give the Tate character of $\Gamma$ the value $u$ at $\gamma$.
For $e\ge1$,
\[
 M(e)_\Gamma\cong
       \mathbb Z_p/(u^{e+1}-1),\qquad
 \operatorname{length}_{\mathbb Z_p}M(e)_\Gamma
       =1+v_p(e+1).
 \tag{HU.5}
\]
For completeness, if $b=1+p^s c$ with $s\ge1$ and $c$ a unit,
the leading term of $b^p-1$ is $p^{s+1}c$ and every later
binomial term has larger valuation.  Here $p$ is odd.
Induction gives $v_p(u^{p^n}-1)=n+1$, and raising such a number
to an exponent prime to $p$ does not change that valuation.
Writing $k=p^nk_0$ proves
$v_p(u^k-1)=1+v_p(k)$ and hence (HU.5).

These finite lengths satisfy
\[
 \begin{split}
 \log p\ \operatorname{length}M(e)_\Gamma
     &\le \log p+\log(e+1),\\
 \sum_{e=1}^{t-1}\operatorname{length}M(e)_\Gamma
     &=(t-1)+v_p(t!)
       \le (t-1)+\frac{t}{p-1}<\frac{3t}{2}.
 \end{split}
 \tag{HU.6}
\]
The factorial estimate follows by summing
$\lfloor t/p^j\rfloor$ over $j\ge1$.
Taking $t=\lfloor x^{401}\rfloor$ therefore gives a contribution
$O(x^{401})$, fully compatible with (HU.3).
Selecting a single $\Delta$-character progression only decreases
the sum.  Allowing polynomially many auxiliary character components
still gives a fixed polynomial in $x$.
Nevertheless
\[
 M/(\gamma^{p^n}-1)M\cong\mathbb Z_p/p^{n+1}\mathbb Z_p,
 \tag{HU.7}
\]
whose lengths tend to infinity.  Reduction gives surjective
transition maps on these quotients, and their inverse limit is $M$.
Thus finite specialization lengths, their polynomial growth and
$\mu=0$ do not force the module to be finite.

\paragraph{Compatibility with a nonzero degree-twelve group.}
This example can respect the relevant parity and exceptional-character
exclusions.  Take the prime
\[
 p=2147483659=2^{31}+11.
\]
Let $\Delta\cong(\mathbb Z/p\mathbb Z)^\times$, write $\omega$
for its Teichm\"uller character, and give $M$ the even
$\Delta$-character $\theta=\omega^{-6}$.
Because $|\Delta|=p-1$ is a $p$-adic unit, the groups
$S_e=(M(e)_\Gamma)^\Delta$ satisfy
\[
 S_e=
 \begin{cases}
  M(e)_\Gamma,&e\equiv6\pmod{p-1},\\
  0,&\text{otherwise}.
 \end{cases}
 \tag{HU.8}
\]
In particular $S_6\cong\mathbb Z/p\mathbb Z$.
The character $\theta$ is neither trivial nor quadratic.
The nonzero positive twists start at $e=6$ and obey
\[
 \log|S_e|\le\log p+\log(e+1)<22+e\le e^7.
\]
Consequently the abstract finite groups $G_k=S_{2k}$ have
\[
 \log|G_k|\le(2k)^7\quad(k\ge1),\qquad G_3\ne0.
 \tag{HU.9}
\]
They satisfy the quantitative bound with constant one and have
prime support above $2^{31}$.  They are a compatibility model
for the numerical statements, not an identification with
$K_{4k}(\mathbb Z)$.

Primality was verified by exact trial division up to $46340$.
An additional Lucas certificate uses
$p-1=2\cdot3\cdot149\cdot2402107$ and witness $2$.
Each factor is prime by trial division,
$2^{p-1}=1\pmod p$, and the residues
$2^{(p-1)/q}\pmod p$, in the displayed factor order, are
\[
 2147483658,\quad2021977949,\quad1547536445,\quad509190093.
\]
Each residue minus one is coprime to $p$.
The exact executable and output are
\path{work/host_strength_model.py} and
\path{work/host_strength_model_results.json}.

\paragraph{Artificial trace polynomials with a nontrivial gcd.}
The two-prime divisibility certificate would become decisive if its
actual evaluated gcd were computed to be one.  Its general numerical
properties do not force this.  To see this directly, for $q=7,13$
put
\[
 E_q=q^{12}+q^6+1,\qquad
 b_q=(-E_q^2\bmod p)\in\{0,\ldots,p-1\},\qquad
 F_q(X)=X^2+b_q.
 \tag{HU.10}
\]
These polynomials are monic integral, their root multisets are
Galois stable, and each root $\alpha$ has
$|\alpha|^2=b_q<p<9q^{12}$.  Hence their normalized roots satisfy
$|\alpha/q^6|<3$.  Their evaluations are nonzero and divisible by
$p$.  The exact data are
\[
 \begin{array}{c|r|r}
 q&E_q&b_q\\ \hline
 7&13841404851&1328819029\\
 13&23298089949291&1446269942
 \end{array}
\]
and
\[
 \begin{split}
 F_7(E_7)&=191584488250615151230,\\
 F_{13}(E_{13})&=542800995285254312397672623,\\
 \gcd(F_7(E_7),F_{13}(E_{13}))
       &=15032385613=7p.
 \end{split}
 \tag{HU.11}
\]
All values and inequalities were checked with exact integers.
The degree-two index sets satisfy the counting upper bounds.
Thus monicity, integrality, Galois stability, the trace
absolute-value bound, nonvanishing, bounded cardinality and
two-prime square divisibility are compatible with a common prime
above the verified range.  These are artificial polynomials:
no claim identifies their roots with automorphic representations
or with a single global compatible system.  Such arithmetic
information is additional to the estimates just tested.

Finally, the available natural-constant bound for $K_{12}$ has
right side $C6^7\ge279936$.  To deduce vanishing from the
verified prime cutoff alone, one would need
$\log|K_{12}|<\log(2^{31})$, or a sufficiently strong
prime-by-prime replacement.  The former bound does not give the
latter.  Thus the remaining gates are actual Iwasawa-module
finiteness, a computed complete trace certificate, or the
separately stated finite-rank homology vanishings.  No new
arithmetic vanishing follows from the Host or window estimates
alone.
